% ═══════════════════════════════════════════════════════════════════
% Paper 0: Γ-Net Framework — Axioms, Quantities, and Verification
% ═══════════════════════════════════════════════════════════════════
\documentclass[aps,pre,twocolumn,superscriptaddress,floatfix]{revtex4-2}
\usepackage{amsmath,amssymb,amsthm}
\usepackage{graphicx}
\usepackage{hyperref}
\usepackage{booktabs}
\usepackage{multirow}
\usepackage{xcolor}

\newtheorem{axiom}{Axiom}
\newtheorem{definition}{Definition}
\newtheorem{theorem}{Theorem}
\newtheorem{proposition}{Proposition}

\begin{document}

\title{$\Gamma$-Net: Minimum Reflection Physics\\
for Dual-Network Life\\[4pt]
\small Framework Definition, Quantity Inventory,\\
and Verification Matrix}

\author{Hsi-Yu Huang}
\email{llc.y.huangll@gmail.com}
\affiliation{Independent Researcher, Taiwan}

\date{\today}

\begin{abstract}
This paper defines the complete axiomatic framework of the
$\Gamma$-Net (\emph{Gamma-Net}) theory: a physics-driven
model of biological organisation based on one variational
principle---Minimum Reflection Action---and three inviolable
constraints C1--C3.
We provide:
(i)~the formal axioms and all derived quantities
($\Gamma$, $A_{\text{imp}}$, $A_{\text{cut}}$,
$D_Z$, $H$, lifecycle ODE, $\Phi_{\text{meta}}$);
(ii)~a unified verification matrix classifying every
prediction as \emph{internal consistency},
\emph{literature-supported}, \emph{NHANES-validated},
or \emph{untested};
(iii)~a falsifiable prediction roadmap for future
experimental work.
The framework unifies phenomena across 15~orders of
magnitude in spatial scale and 19~tissue types under
a single equation: $\Delta Z = -\eta\,\Gamma\,x_{\text{in}}\,x_{\text{out}}$.
This paper contains no new results; it serves as the
master reference for the five companion papers that
develop and test the theory.
\end{abstract}

\maketitle

% ============================================================
%  SECTION 1 — INTRODUCTION
% ============================================================
\section{Introduction: Why This Paper Exists}
\label{sec:intro}

The $\Gamma$-Net framework has been developed across five
companion papers, each deriving a distinct class of
biological phenomena from the Minimum Reflection Principle
(MRP).
Paper~1 derives topology and metabolic scaling~\cite{Paper1};
Paper~2 derives dual neural--vascular organ
health~\cite{Paper2};
Paper~3 derives sleep, the lifecycle equation, and
aging~\cite{Paper3};
Paper~4 derives memory, consciousness, and
nociception~\cite{Paper4};
Paper~5 validates the framework against NHANES
epidemiological data and 19~tissue remodeling
laws~\cite{Paper5}.

A reader entering at any single paper encounters
a local recap of the axioms but lacks:
\begin{enumerate}
  \item a single, canonical list of \emph{all} defined
    quantities and their inter-relations;
  \item a structured table distinguishing which claims
    have external (non-circular) support and which
    remain internal consistency checks;
  \item a roadmap of falsifiable predictions organised
    by experimental modality.
\end{enumerate}
This paper fills those three gaps.
It contains no new derivations or simulations; every result
cited here is proven elsewhere and cross-referenced
to the relevant companion paper.


% ============================================================
%  SECTION 2 — AXIOMS
% ============================================================
\section{The Three Axioms}
\label{sec:axioms}

The theory rests on one variational principle and three
local constraints.

\subsection{The Minimum Reflection Principle (MRP)}
\label{sec:mrp}

\begin{axiom}[Minimum Reflection Action]
\label{ax:mrp}
Any impedance network subject to sustained energy flow
evolves so as to minimise the total reflected action:
\begin{equation}
  \mathcal{A}[\Gamma]
  = \int_0^T \sum_i \Gamma_i^2(t)\,P_{\text{in},i}(t)\,dt
  \;\to\;\min\,,
  \label{eq:action}
\end{equation}
where $\Gamma_i = (Z_{\text{load},i} - Z_{\text{source},i})
/ (Z_{\text{load},i} + Z_{\text{source},i})$ is the
reflection coefficient at interface~$i$, and
$P_{\text{in},i}$ is the incident power.
\end{axiom}

\noindent
The MRP is not a modelling choice; it is the unique
variational formulation consistent with the First Law of
Thermodynamics at a one-dimensional interface
(Paper~5, Section~2~\cite{Paper5}).

\subsection{C1 --- Energy Conservation}
\label{sec:c1}

\begin{axiom}[C1: Energy Conservation]
\label{ax:c1}
At every interface $i$, at every instant $t$:
\begin{equation}
  \Gamma_i^2(t) + T_i(t) = 1\,,
  \label{eq:c1}
\end{equation}
where $T_i = 1 - \Gamma_i^2$ is the transmission
coefficient (fraction of incident power delivered
to the load).
\end{axiom}

\noindent
C1 is the First Law of Thermodynamics applied locally.
It is not an approximation; it holds exactly for any
lossless interface.  Dissipative corrections (Johnson--Nyquist
noise, Arrhenius aging) modify $Z$ over time but never
violate power balance at a single interface.

\subsection{C2 --- Impedance Remodeling}
\label{sec:c2}

\begin{axiom}[C2: Impedance Remodeling]
\label{ax:c2}
The gradient descent of $\mathcal{A}$ with respect to
local impedance yields:
\begin{equation}
  \Delta Z_i = -\eta\,\Gamma_i\,x_{\text{in},i}\,x_{\text{out},i}\,,
  \label{eq:c2}
\end{equation}
where $\eta > 0$ is the remodeling rate,
$x_{\text{in},i}$ is the input signal (effort or flow
arriving at interface~$i$), and
$x_{\text{out},i}$ is the output signal (transmitted
through interface~$i$).
\end{axiom}

\noindent
Equation~\eqref{eq:c2} is the unique local update rule
that decreases $\mathcal{A}$ at every step
(Paper~5, Section~2~\cite{Paper5}).
Its form is identical to Hebb's learning rule when $Z$
is synaptic conductance, to Wolff's law when $Z$ is bone
stiffness, to Glagov's remodeling when $Z$ is vascular
compliance, and to 16~further tissue-specific remodeling
laws (Table~\ref{tab:nineteen}).

\medskip\noindent
\textbf{Historical note.}
In the original formulation (Papers~1--4), C2 was called
the ``Hebbian update'' and its signals were labelled
$x_{\text{pre}}$, $x_{\text{post}}$.
Paper~5 proved that the same equation governs all 
19~tissue types examined, making the neural-specific
terminology misleading.
We adopt the generalised names
\emph{impedance remodeling} and
$(x_{\text{in}}, x_{\text{out}})$ throughout this series.

\subsection{C3 --- Impedance-Tagged Transport}
\label{sec:c3}

\begin{axiom}[C3: Impedance-Tagged Transport]
\label{ax:c3}
Every energy or material flow across an interface carries
$(Z_{\text{source}}, Z_{\text{load}})$ as metadata,
so that $\Gamma_i$ is locally computable without
global network knowledge.
\end{axiom}

\noindent
C3 is a bookkeeping requirement, not an additional
physical law.  It states that whenever a signal crosses
an impedance boundary, the downstream node can compute
$\Gamma$ from local information alone.
In biological terms: blood carries its viscosity
($Z_{\text{blood}}$) into a vessel of compliance
($Z_{\text{vessel}}$); a nerve impulse carries
$R_m + 1/j\omega C_m$ into a synapse of input impedance
$Z_{\text{syn}}$; a mechanical load carries force/velocity
into a bone of stiffness $Z_{\text{bone}}$.

\medskip\noindent
\textbf{Implementation note.}
In the open-source reference implementation (ALICE),
C3 is enforced by requiring all inter-module values to be
\texttt{ElectricalSignal} objects carrying $Z$ metadata.
Raw floating-point numbers between modules are forbidden
at the code level.


% ============================================================
%  SECTION 3 — QUANTITY INVENTORY
% ============================================================
\section{Derived Quantities}
\label{sec:quantities}

Table~\ref{tab:quantities} lists every named quantity
in the $\Gamma$-Net framework, its defining equation,
physical interpretation, and the companion paper where
it first appears.

\begin{table*}[t]
\caption{Complete quantity inventory of the $\Gamma$-Net framework.}
\label{tab:quantities}
\begin{ruledtabular}
\begin{tabular}{llllc}
\textbf{Symbol} & \textbf{Name} & \textbf{Definition} & \textbf{Interpretation} & \textbf{Paper} \\
\hline
$\Gamma_i$ & Reflection coefficient
  & $\frac{Z_{\text{load},i}-Z_{\text{source},i}}{Z_{\text{load},i}+Z_{\text{source},i}}$
  & Mismatch at interface $i$ & 0 \\[4pt]
%
$T_i$ & Transmission coefficient
  & $1 - \Gamma_i^2$
  & Fraction of power delivered & 0 \\[4pt]
%
$\mathcal{A}$ & Reflection action
  & $\int_0^T\!\sum_i \Gamma_i^2 P_{\text{in},i}\,dt$
  & Total wasted energy & 0 \\[4pt]
%
$A_{\text{imp}}$ & Improvable action
  & $\mathcal{A} - A_{\text{cut}}$
  & Reducible by C2 remodeling & 1 \\[4pt]
%
$A_{\text{cut}}$ & Irreducible (cutoff) action
  & $A_{\text{cut}} = \sum_{k>K^*}\!\Gamma_k^2$
  & Topological cost immune to C2 & 1 \\[4pt]
%
$D_Z$ & Impedance debt
  & $\int_0^{T_{\text{wake}}}\!|\Gamma|^2 P_{\text{in}}\,dt$
  & Cumulative mismatch; cleared by sleep & 3 \\[4pt]
%
$D_n,\,D_v$ & Neural / vascular debt
  & $D_Z = D_n + D_v$
  & Dual-network decomposition & 3 \\[4pt]
%
$H$ & Health index
  & $\prod_i(1-\Gamma_i^2)$
  & Whole-body transmission efficiency & 5 \\[4pt]
%
$\Sigma\Gamma^2$ & Total mismatch energy
  & $\sum_i \Gamma_i^2(t)$
  & Instantaneous network mismatch & 3 \\[4pt]
%
\multirow{2}{*}{Lifecycle ODE} & \multirow{2}{*}{Lifecycle equation}
  & $\frac{d(\Sigma\Gamma^2)}{dt} = -\eta\,\Sigma\Gamma^2$
  & \multirow{2}{*}{Cognition bathtub curve} & \multirow{2}{*}{3} \\
  & & $\;+\gamma\,\Gamma_{\text{env}} + \delta\,D(t)$  & & \\[4pt]
%
$\Phi_{\text{meta}}$ & Meta-reflection coefficient
  & $\frac{Z_{\text{obs}}-Z_{\text{act}}}{Z_{\text{obs}}+Z_{\text{act}}}$
  & Self-monitoring; consciousness marker & 4 \\[4pt]
%
$\mathcal{S}$ & Null space (``soul'')
  & $\ker(\partial Z / \partial t)$
  & Impedances immune to all experience & 4 \\[4pt]
%
$\beta$ & Murray bifurcation ratio
  & $2^{1/3} \approx 1.26$
  & Optimal branch radius ratio & 1 \\[4pt]
%
$\Omega$ & Dual-network overhead
  & $\ge 1$
  & Packing cost of two trees in one volume & 1 \\
\end{tabular}
\end{ruledtabular}
\end{table*}


% ============================================================
%  SECTION 4 — SERIES MAP
% ============================================================
\section{Series Architecture}
\label{sec:series}

The five companion papers form a logical chain:

\begin{enumerate}
  \item \textbf{Paper~1: Topology and Fractal Cost}~\cite{Paper1}.
    Proves the irreducibility theorem
    ($\mathcal{A} = A_{\text{imp}} + A_{\text{cut}}$,
    $A_{\text{cut}} > 0$), derives relay-node hierarchies,
    the Murray bifurcation ratio $\beta = 2^{1/3}$,
    Kleiber's 3/4-power metabolic scaling, joint impedance
    transformers, and dual-network packing overhead~$\Omega$.

  \item \textbf{Paper~2: Dual-Network Organ Health}~\cite{Paper2}.
    Applies MRP to vascular networks; derives Murray's law
    from $\Gamma$-minimisation; defines the dual-network
    health index $H = (1-\Gamma_n^2)(1-\Gamma_v^2)$;
    models five vascular diseases as impedance failure modes.
    Validated against 12/12 published clinical benchmarks with
    100\% pass rate.

  \item \textbf{Paper~3: Temporal Dynamics}~\cite{Paper3}.
    Introduces impedance debt $D_Z$; derives sleep as
    $D_Z$-discharge; derives the lifecycle equation governing
    cognition from birth to death; models emotion as
    $\Gamma$-dynamics readout; derives aging as Arrhenius
    fatigue; explains aging pain immunity via slow $Z$-drift.

  \item \textbf{Paper~4: Memory, Consciousness, Nociception}~\cite{Paper4}.
    Defines memory as $dZ/dt$, consciousness as closed-loop
    $\Phi_{\text{meta}}$, and the null space $\mathcal{S}$
    (``soul'') as $\ker(\partial Z/\partial t)$.
    Proves the Television Theorem distinguishing open-loop
    from closed-loop systems.

  \item \textbf{Paper~5: Clinical Validation}~\cite{Paper5}.
    Proves C2 governs all 19~tissue types (not just neurons);
    validates against NHANES ($n = 49{,}774$, zero fitted
    parameters, 7/7~organs significant);
    reports all-cause mortality $C = 0.70$,
    Q1-vs-Q4 risk ratio $= 6.89\times$;
    provides sex-as-boundary-condition analysis.
\end{enumerate}

\noindent
Papers~1--4 are primarily theoretical with internal
computational verification.
Paper~5 breaks the circularity with external epidemiological
data.


% ============================================================
%  SECTION 5 — VERIFICATION MATRIX
% ============================================================
\section{Verification and Validation Matrix}
\label{sec:vv}

Table~\ref{tab:vv} classifies every major claim by its
current evidence status.
Four levels are used:

\begin{description}
  \item[IC] \textbf{Internal consistency.}
    Verified within the ALICE simulation engine.
    Self-confirming; does not constitute independent evidence.
  \item[LIT] \textbf{Literature-supported.}
    Prediction matches published experimental or clinical
    data not used in model construction.
  \item[EXT] \textbf{Externally validated.}
    Tested against independent dataset (NHANES) with
    zero fitted parameters and hash-locked predictions.
  \item[UT] \textbf{Untested.}
    Falsifiable prediction not yet confronted with data.
\end{description}

\begin{table*}[t]
\caption{Verification and validation matrix.
IC = internal consistency;
LIT = literature-supported;
EXT = externally validated (NHANES);
UT = untested prediction.}
\label{tab:vv}
\begin{ruledtabular}
\begin{tabular}{llcccc}
\textbf{Claim} & \textbf{Paper} & \textbf{IC} & \textbf{LIT} & \textbf{EXT} & \textbf{UT} \\
\hline
\multicolumn{6}{l}{\textit{Topology and scaling (Paper~1)}} \\
Relay nodes emerge from MRP & 1 & $\checkmark$ & & & \\
$\tau_{\text{conv}} \propto N^{-0.46}$ & 1 & $\checkmark$ & & & \\
$A_{\text{cut}} > 0$ (irreducibility) & 1 & $\checkmark$ & & & \\
Murray ratio $\beta = 2^{1/3}$ & 1 & $\checkmark$ & $\checkmark$ & & \\
Kleiber's law $B \propto M^{3/4}$ & 1 & $\checkmark$ & $\checkmark$ & & \\
Joint = quarter-wave transformer & 1 & $\checkmark$ & & & $\checkmark$ \\
$A_{\text{cut}} \propto M^{2/3}$ & 1 & $\checkmark$ & $\checkmark$ & & \\
\hline
\multicolumn{6}{l}{\textit{Dual-network organ health (Paper~2)}} \\
Murray's law from $\Gamma$-minimisation & 2 & $\checkmark$ & $\checkmark$ & & \\
Dual-network health $H=(1{-}\Gamma_n^2)(1{-}\Gamma_v^2)$ & 2 & $\checkmark$ & & $\checkmark$ & \\
Vascular disease as $Z$-mismatch & 2 & $\checkmark$ & $\checkmark$ & & \\
12/12 clinical calibration tests & 2 & & $\checkmark$ & & \\
\hline
\multicolumn{6}{l}{\textit{Temporal dynamics (Paper~3)}} \\
Sleep = $D_Z$ discharge & 3 & $\checkmark$ & $\checkmark$ & & \\
Jellyfish sleep without brain & 3 & & $\checkmark$ & & \\
Brain evolved for thermal management & 3 & $\checkmark$ & $\checkmark$ & & \\
Lifecycle bathtub curve & 3 & $\checkmark$ & $\checkmark$ & & \\
Yerkes--Dodson as impedance optimum & 3 & $\checkmark$ & $\checkmark$ & & \\
Aging = Arrhenius fatigue on $Z$ & 3 & $\checkmark$ & & & \\
Aging pain immunity (slow drift) & 3 & $\checkmark$ & & & $\checkmark$ \\
Infant sleep duration vs $D_Z$ & 3 & $\checkmark$ & $\checkmark$ & & \\
\hline
\multicolumn{6}{l}{\textit{Memory, consciousness, nociception (Paper~4)}} \\
Memory = $dZ/dt$ & 4 & $\checkmark$ & & & \\
$\Phi_{\text{meta}}$ as consciousness marker & 4 & $\checkmark$ & & & $\checkmark$ \\
Null space $\mathcal{S} = \ker(\partial Z/\partial t)$ & 4 & $\checkmark$ & & & \\
Television Theorem (open/closed loop) & 4 & $\checkmark$ & & & $\checkmark$ \\
Thermal hierarchy: PFC fails before V1 & 4 & $\checkmark$ & $\checkmark$ & & $\checkmark$ \\
\hline
\multicolumn{6}{l}{\textit{Clinical validation (Paper~5)}} \\
C2 = 19 tissue remodeling laws & 5 & & $\checkmark$ & & \\
NHANES 7/7 organ AUC significant & 5 & & & $\checkmark$ & \\
Endocrine AUC = 0.823 & 5 & & & $\checkmark$ & \\
All-cause mortality $C = 0.70$ & 5 & & & $\checkmark$ & \\
Q1 vs Q4 mortality ratio = $6.89\times$ & 5 & & & $\checkmark$ & \\
Sex as boundary condition on $Z_{\text{setpoint}}$ & 5 & & $\checkmark$ & $\checkmark$ & \\
Cardiac AUC gain $+0.166$ from vascular input & 5 & & & $\checkmark$ & \\
$H$ decreases with disease burden ($\rho = -0.378$) & 5 & & & $\checkmark$ & \\
\end{tabular}
\end{ruledtabular}
\end{table*}


% ============================================================
%  SECTION 6 — PREDICTION ROADMAP
% ============================================================
\section{Prediction Roadmap}
\label{sec:roadmap}

Table~\ref{tab:predictions} organises the falsifiable
predictions of the $\Gamma$-Net framework by experimental
modality, enabling independent groups to test specific
claims without needing the full simulation engine.

\begin{table*}[t]
\caption{Falsifiable prediction roadmap.
Each prediction is testable with existing technology.}
\label{tab:predictions}
\begin{ruledtabular}
\begin{tabular}{lllc}
\textbf{Domain} & \textbf{Prediction} & \textbf{Method} & \textbf{Paper} \\
\hline
\multicolumn{4}{l}{\textit{Physiological}} \\
& PFC fails before V1 under hyperthermia
  & Controlled heating + fMRI/EEG & 4 \\
& Vascular $\Gamma$ correlates with PWV / elastography
  & BIS + arterial tonometry & 2 \\
& Joint cartilage $\approx$ quarter-wave transformer
  & Ultrasound impedance spectroscopy & 1 \\
& Bone-on-bone OA raises local $\Gamma^2$ measurably
  & Intra-articular force sensors & 1 \\
\hline
\multicolumn{4}{l}{\textit{Sleep \& circadian}} \\
& Cross-species $D_{\text{thermal}}/D_{\text{micro}}$ ratio
  predicts NREM/REM partition
  & Polysomnography + metabolic calorimetry & 3 \\
& Sleep rebound magnitude $\propto$ wake-duration $D_Z$
  & EEG slow-wave power after controlled deprivation & 3 \\
& Jellyfish sleep duration scales with $\kappa_{\text{eff}}$
  & Temperature-controlled aquarium experiments & 3 \\
\hline
\multicolumn{4}{l}{\textit{Cognitive / developmental}} \\
& Infantile amnesia offset tracks myelination Z-change
  & Retrospective longitudinal EEG + structural MRI & 3 \\
& Yerkes--Dodson peak shifts with $\eta$ (skill level)
  & Psychophysics + cortisol monitoring & 3 \\
& Aging pain immunity: slow $Z$-drift below pain threshold
  & Quantitative sensory testing, longitudinal & 3 \\
\hline
\multicolumn{4}{l}{\textit{Clinical / epidemiological}} \\
& $\Gamma$-vector predicts disease in non-US populations
  & NHANES-equivalent cohorts (UK Biobank, Taiwan MJ) & 5 \\
& $H$ predicts surgical outcome / length of stay
  & Prospective peri-operative cohort & 5 \\
& Network propagation: renal $\Gamma$ predicts cardiac events
  & Longitudinal cardiorenal registry & 5 \\
\end{tabular}
\end{ruledtabular}
\end{table*}


% ============================================================
%  SECTION 7 — INTERFACE SCALE TABLE (SUMMARY)
% ============================================================
\section{The Interface Scale Table}
\label{sec:scales}

The same equation ($\Gamma$, C1, C2) applies across seven
distinct physical domains and 14~orders of magnitude in
spatial scale:

\begin{table}[h]
\caption{Representative interfaces grouped by physical domain.
Full table with impedance definitions: Paper~5~\cite{Paper5}.}
\label{tab:scales}
\begin{ruledtabular}
\begin{tabular}{llc}
\textbf{Domain} & \textbf{Interface} & \textbf{Scale} \\
\hline
Electrical   & Synapse / ion channel       & 25\,nm--1\,$\mu$m \\
Electrical   & Power-grid transformer      & 10\,m \\
Hydraulic    & Arteriole (Poiseuille)      & 100\,$\mu$m \\
Hydraulic    & Whole-body vascular         & 1\,m \\
Acoustic     & Tissue boundary (ultrasound)& 1\,mm \\
Acoustic     & Submarine hull              & 10\,m \\
Electromagnetic & Optical fibre splice     & 100\,$\mu$m \\
Electromagnetic & Ionospheric waveguide    & $10^6$\,m \\
Mechanical   & Bone--cartilage joint       & 10\,cm \\
Mechanical   & Tendon--muscle junction     & 1\,cm \\
Thermal      & Brain thermal boundary      & 10\,cm \\
Thermal      & Skin--air interface         & 1\,m \\
Chemical     & Organ parenchyma            & 10\,cm \\
\end{tabular}
\end{ruledtabular}
\end{table}

\noindent
At every scale, the physics is identical: only the
definition of $Z$ changes (electrical, acoustic,
mechanical, hydraulic, thermal, electromagnetic,
chemical).
The universality is not an analogy; it is the same
theorem applied to different media.

Moreover, every row in the table participates in the same
closed loop: \textbf{perception}
($x_{\text{in}}$ measurement) $\to$
\textbf{error} ($\Gamma$ mismatch) $\to$
\textbf{compensation} (C2: $\Delta Z =
-\eta\,\Gamma\,x_{\text{in}}\,x_{\text{out}}$) $\to$
\textbf{re-perception} (updated $\Gamma'$ and
$x_{\text{in}}'$).
This perception--compensation cycle is the operational
meaning of C1--C2 at every interface; Paper~5
elaborates it across all seven domains.


% ============================================================
%  SECTION 8 — NINETEEN TISSUES (SUMMARY TABLE)
% ============================================================
\section{The Nineteen C2 Networks}
\label{sec:nineteen}

Paper~5 proves that C2 (Eq.~\ref{eq:c2}) recovers 19
independently discovered tissue-remodeling laws when the
generic signals $(x_{\text{in}}, x_{\text{out}})$ are
identified with tissue-specific physical quantities.
Table~\ref{tab:nineteen} provides the mapping.

\begin{table*}[t]
\caption{The 19 $\Gamma$-networks.
Each row maps C2 to a known remodeling law by identifying
$x_{\text{in}}$, $x_{\text{out}}$, and $Z$.
Full derivations: Paper~5~\cite{Paper5}.}
\label{tab:nineteen}
\begin{ruledtabular}
\begin{tabular}{lllll}
\textbf{Tissue} & \textbf{Known law}
  & $x_{\text{in}}$ & $x_{\text{out}}$ & $Z$ \textbf{definition} \\
\hline
Neural        & Hebb's law (1949)              & Pre-synaptic spike & Post-synaptic EPSP & $R_m + 1/j\omega C_m$ \\
Vascular      & Glagov remodeling (1987)       & Wall shear stress  & Endothelial response & $P/Q$ (Poiseuille) \\
Skeletal      & Wolff's law (1892)             & Applied load       & Bone strain        & Force/velocity \\
Muscular      & Davis's law (1867)             & Neural drive       & Muscle tension     & Stiffness \\
Immune        & Clonal selection               & Antigen (APC)      & T/B cell activation & Receptor affinity \\
Epithelial    & Keratinisation                 & External friction  & Tissue deformation & Thermal+mechanical $Z$ \\
Pulmonary     & Altitude acclimatisation       & Ventilation demand & Alveolar expansion & $P_{\text{aw}}/\dot{V}$ \\
GI tract      & Intestinal adaptation          & Luminal content    & Epithelial uptake  & Absorption $Z$ \\
Endocrine     & Receptor regulation            & Hormone conc.      & Intracellular resp.& $1/K_d$ \\
Lymphatic     & Reactive hypertrophy           & Tissue load        & Lymph node response& Drainage $Z$ \\
Renal         & TG feedback                    & Macula densa signal& Arteriolar response& $P_{\text{glom}}/Q_{\text{filt}}$ \\
Hepatic       & Liver regeneration             & Portal HGF         & Hepatocyte prolif. & Clearance $Z$ \\
Connective    & Tendon remodeling              & Mechanical strain  & Fibroblast response& $EA/L$ \\
Adipose       & Leptin signaling               & Insulin signal     & Adipocyte response & Storage $Z$ \\
Hematopoietic & Erythropoiesis                 & EPO signal         & Erythroid response & $P_{O_2}/Q_{O_2}$ \\
Cartilage     & \textit{C2 null} ($\eta\approx 0$) & Mechanical load & (negligible)       & $H_A A/h$ \\
Reproductive  & Follicular selection           & Gonadotropin       & Granulosa/germ cell& $1/(N_R K_a)$ \\
Cardiac       & Laplace--Starling              & Wall stress $\sigma$ & Myocardial strain $\varepsilon$ & Chamber compliance \\
Brain (reg.)  & Baroreflex resetting           & Sensory afferent   & Motor efferent     & Controller gain \\
\end{tabular}
\end{ruledtabular}
\end{table*}


% ============================================================
%  SECTION 9 — NOTATION
% ============================================================
\section{Notation Conventions}
\label{sec:notation}

\begin{itemize}
  \item $Z$: impedance (effort/flow ratio). Units depend on
    medium: $\Omega$ (electrical), Pa$\cdot$s/m$^3$ (vascular),
    N$\cdot$s/m (mechanical).
  \item $\Gamma$: dimensionless, $\in [-1, +1]$.
  \item $\eta$: remodeling rate, $> 0$. Tissue-dependent;
    $\eta \approx 0$ for cartilage (C2 null case).
  \item $(x_{\text{in}}, x_{\text{out}})$: generic input and
    output signals at an interface.  Replaces the earlier
    $(x_{\text{pre}}, x_{\text{post}})$ notation.
  \item $D_Z$: impedance debt, units of energy $\times$ time
    (J$\cdot$s).
  \item $H$: dimensionless, $\in [0, 1]$. Higher is healthier.
\end{itemize}


% ============================================================
%  SECTION 10 — CONCLUSION
% ============================================================
\section{Conclusion}
\label{sec:conclusion}

The $\Gamma$-Net framework rests on one variational
principle (MRP) and three local constraints
(C1: energy conservation; C2: impedance remodeling;
C3: impedance-tagged transport).
From these, five companion papers derive topology,
organ health, temporal dynamics, consciousness, and
clinical diagnostics---spanning 15 orders of spatial
magnitude and 19 tissue types.

This paper provides the master reference:
a canonical quantity list (Table~\ref{tab:quantities}),
a verification matrix distinguishing internal from
external evidence (Table~\ref{tab:vv}),
and a falsifiable prediction roadmap
(Table~\ref{tab:predictions}).

All source code, simulation scripts, and NHANES pipelines
are available at
\url{https://github.com/cyhuang76/alice-gamma-net}
under AGPL-3.0.
The full test suite comprises 3{,}274 automated tests
with 100\% pass rate.


% ============================================================
%  BIBLIOGRAPHY
% ============================================================
\begin{thebibliography}{99}

\bibitem{Paper1}
H.-Y.~Huang,
``Topology and fractal cost of $\Gamma$-networks:
relay nodes, metabolic scaling, and the irreducibility
theorem,''
companion Paper~1 (2025).

\bibitem{Paper2}
H.-Y.~Huang,
``Dual-network organ health: vascular $\Gamma$-physics
and the five failure modes,''
companion Paper~2 (2025).

\bibitem{Paper3}
H.-Y.~Huang,
``Temporal dynamics of impedance debt: sleep, the lifecycle
equation, and the physics of aging,''
companion Paper~3 (2025).

\bibitem{Paper4}
H.-Y.~Huang,
``Memory, consciousness, and nociception as impedance
dynamics,''
companion Paper~4 (2025).

\bibitem{Paper5}
H.-Y.~Huang,
``Clinical validation: nineteen tissue remodeling laws,
NHANES epidemiology, and mortality prediction,''
companion Paper~5 (2025).

\bibitem{ALICE_code}
H.-Y.~Huang,
\emph{ALICE: A physics-driven electronic lifeform},
\url{https://github.com/cyhuang76/alice-gamma-net} (2025).

\end{thebibliography}

\end{document}
